\chapter{Conclusion}
\label{chap-conclusion}

La première version de ce projet répondait à la question posée en 2023 : savons-nous
construire une interface graphique avec GTK ? Elle y répondait bien --- le jeu était
complet --- mais elle laissait de côté l'autre moitié du cours, celle qui parle
d'organisation : un concept par module, une utilisation indépendante de la réalisation, une
fonction écrite donc testée.

La seconde version répond à cette autre moitié. Ce qu'elle démontre tient en trois points.

\textbf{Séparer coûte peu et rapporte beaucoup.} Extraire les règles dans une bibliothèque
sans dépendance a rendu possibles, d'un seul coup, les tests automatiques, la seconde
interface, la mesure des stratégies et la reproductibilité. Aucune de ces quatre choses
n'était envisageable tant qu'une case du jeu contenait l'adresse d'un bouton.

\textbf{Mesurer corrige l'intuition.} Nous pensions, en 2023, que la force de l'ordinateur
tenait à sa façon de poursuivre un bateau touché. Les mesures montrent que l'essentiel est
ailleurs : respecter ce que la règle de non-contact prouve fait gagner trente tirs, la
poursuite moins de huit. Les niveaux de difficulté ont été redéfinis d'après ces chiffres.

\textbf{Les outils trouvent ce que l'usage cache.} La lecture hors limites du niveau
difficile n'a jamais fait planter le jeu devant nous. Elle était pourtant là, et un outil
standard l'aurait signalée dès le premier lancement. Vérifier automatiquement, à chaque
modification, n'est pas un supplément de confort : c'est ce qui permet de modifier un
programme sans craindre de le casser.

Le jeu, lui, n'a pas changé. On y commande toujours la marine de Sup Galilée contre les
pirates, sur une grille de dix cases sur dix. C'était le but : que le joueur ne voie pas la
différence, et que le prochain développeur ne voie qu'elle.
